% =========================================================
\chapter{Fundamentação Teórica}
% =========================================================

A Fundamentação Teórica deve apresentar os conhecimentos indispensáveis para
que o leitor compreenda o problema, a solução proposta e os resultados
posteriormente apresentados.

% ---------------------------------------------------------
\section{Considerações iniciais}
% ---------------------------------------------------------

\begin{enumerate}
	\item Quais conceitos são indispensáveis para compreender o trabalho?
	\item Quais conhecimentos são pressupostos?
	\item Quais conceitos precisam ser formalmente definidos?
	\item Quais teorias sustentam a solução proposta?
	\item Quais variáveis e grandezas serão utilizadas posteriormente?
	\item Quais conceitos serão utilizados no modelo matemático?
	\item Quais conceitos serão utilizados nas simulações?
	\item Qual é a relação entre os fundamentos teóricos apresentados?
	\item O capítulo está fornecendo apenas o conhecimento necessário ou está se transformando em um tratado sobre o assunto?
	\item Como a fundamentação prepara o leitor para compreender a solução?
\end{enumerate}

% ---------------------------------------------------------
\section{Conteúdo seminal do assunto correlato}
% ---------------------------------------------------------

\begin{enumerate}
	\item Qual é a origem histórica do conceito ou tecnologia?
	\item Quem foram os pesquisadores responsáveis pelas contribuições fundamentais?
	\item Quais trabalhos são considerados seminais?
	\item Qual problema originalmente motivou o desenvolvimento da teoria?
	\item Quais são os princípios físicos, matemáticos ou computacionais envolvidos?
	\item Quais equações fundamentais descrevem o fenômeno?
	\item Quais hipóteses e simplificações estão presentes nos modelos?
	\item Quais são as limitações da teoria?
	\item Como a teoria evoluiu até o estado atual?
	\item Quais fundamentos serão efetivamente utilizados na solução proposta?
\end{enumerate}

{Observação.}
O termo ``seminal'' não significa necessariamente ``antigo''. O critério deve
ser a relevância fundamental do conteúdo para a compreensão do problema, e
não simplesmente sua antiguidade.

% ---------------------------------------------------------
\section{Considerações finais}
% ---------------------------------------------------------

\begin{enumerate}
	\item Quais fundamentos teóricos são essenciais para o trabalho?
	\item Quais equações serão utilizadas posteriormente?
	\item Quais conceitos serão utilizados no modelo?
	\item Quais parâmetros deverão ser definidos experimentalmente ou por simulação?
	\item Quais hipóteses teóricas serão adotadas?
	\item Quais limitações teóricas podem afetar os resultados?
	\item Como a fundamentação sustenta a solução proposta?
	\item O leitor possui, ao final deste capítulo, conhecimento suficiente para compreender o próximo?
\end{enumerate}